\documentclass{article}
\usepackage[utf8]{inputenc}
\usepackage{amsmath}
\usepackage{amssymb}
\usepackage{enumitem}

\title{Proof Export}
\date{}

\begin{document}
\maketitle

\section*{Node 1}

\textbf{Statement:} (Barvinok-Pataki for Order Unit Spaces). Let A be a unital *-algebra (or more generally an Archimedean order unit space (V, V\textasciicircum{}+, e)) and let a\_1, ..., a\_m ∈ A\textasciicircum{}\{sa\} be self-adjoint elements. Define the affinely constrained state space C = \{φ ∈ S(A) : φ(a\_j) = β\_j, j = 1, ..., m\} (or inequality constraints φ(a\_j) ≤ β\_j, or mixed). If φ is an extreme point of C, then:

(I) dim\_R(π\_φ(A)'\textasciicircum{}\{sa\}) ≤ m + 1,

where π\_φ(A)' is the commutant of the GNS representation associated to φ, and (·)\textasciicircum{}\{sa\} denotes the real subspace of self-adjoint elements.

(II) Equivalently, if π\_φ ≅ ⊕\_\{j=1\}\textasciicircum{}k σ\_j\textasciicircum{}\{⊕n\_j\} is the decomposition into pairwise inequivalent irreducible representations σ\_j with multiplicities n\_j, then:

Σ\_\{j=1\}\textasciicircum{}k n\_j² ≤ m + 1.

Moreover, this result simultaneously recovers: (a) Carathéodory's theorem (abelian case A = C(X)), (b) the classical Barvinok-Pataki bound for complex SDP (A = M\_n(C)), (c) a new mixed Carathéodory-Barvinok-Pataki bound for block-diagonal algebras (A = ⊕ M\_\{d\_i\}(C)), and (d) the Weis-Shirokov theorem on extreme points of energy-constrained quantum states (A = B(H)).

\textbf{Type:} claim

\textbf{Inference:} assumption

\textbf{Status:} pending

\textbf{Taint:} unresolved

\begin{enumerate}
\item \textbf{Node 1.1:} (Face-Commutant Correspondence — Alfsen-Shultz). Let A be a unital *-algebra and φ ∈ S(A) a state with GNS triple (π\_φ, H\_φ, ξ\_φ). Then the face F\_\{S(A)\}(φ) generated by φ in the state space S(A) is affinely isomorphic to the set \{T ∈ π\_φ(A)'\textasciicircum{}+ : ⟨Tξ\_φ, ξ\_φ⟩ = 1\} via the map T ↦ ω\_T defined by ω\_T(a) = ⟨π\_φ(a)Tξ\_φ, ξ\_φ⟩. Consequently, dim(F\_\{S(A)\}(φ)) = dim\_R(π\_φ(A)'\textasciicircum{}\{sa\}) - 1.

  Type: claim, Inference: assumption, Status: pending, Taint: unresolved

\begin{enumerate}
\item \textbf{Node 1.1.1:} (Well-definedness of the map — 2nd amendment, cleaned). The map Φ : \{T ∈ π\_φ(A)'\textasciicircum{}+ : ⟨Tξ\_φ, ξ\_φ⟩ = 1\} → S(A) defined by Φ(T) = ω\_T where ω\_T(a) = ⟨π\_φ(a)Tξ\_φ, ξ\_φ⟩ is well-defined: ω\_T is a state on A.

Proof: (i) Linearity in a: clear from linearity of π\_φ and linearity of the inner product in its first argument.

(ii) Positivity: ω\_T(a*a) = ⟨π\_φ(a*a)Tξ\_φ, ξ\_φ⟩ = ⟨π\_φ(a)*π\_φ(a)Tξ\_φ, ξ\_φ⟩ (since π\_φ is a *-homomorphism). Since T ∈ π\_φ(A)', we have Tπ\_φ(a) = π\_φ(a)T, hence π\_φ(a)T = Tπ\_φ(a), giving ⟨π\_φ(a)*Tπ\_φ(a)ξ\_φ, ξ\_φ⟩ = ⟨Tπ\_φ(a)ξ\_φ, π\_φ(a)ξ\_φ⟩ ≥ 0, where the last step uses T ≥ 0 (⟨Tη, η⟩ ≥ 0 for all η, with η = π\_φ(a)ξ\_φ).

(iii) Normalisation: ω\_T(1) = ⟨π\_φ(1)Tξ\_φ, ξ\_φ⟩ = ⟨Tξ\_φ, ξ\_φ⟩ = 1 by the domain constraint. This is a direct proof, not an assumption.

  Type: claim, Inference: assumption, Status: pending, Taint: unresolved

\item \textbf{Node 1.1.2:} (The map Φ is affine — amended). The domain D = \{T ∈ π\_φ(A)'\textasciicircum{}+ : ⟨Tξ\_φ, ξ\_φ⟩ = 1\} is convex: for T\_1, T\_2 ∈ D and 0 ≤ λ ≤ 1, let S = λT\_1 + (1-λ)T\_2. Then: (a) S ∈ π\_φ(A)' since the commutant is a linear subspace (von Neumann algebra); (b) S ≥ 0 since λ, (1-λ) ≥ 0 and T\_1, T\_2 ≥ 0; (c) ⟨Sξ\_φ, ξ\_φ⟩ = λ⟨T\_1ξ\_φ, ξ\_φ⟩ + (1-λ)⟨T\_2ξ\_φ, ξ\_φ⟩ = λ·1 + (1-λ)·1 = 1 (using linearity of the inner product in the first argument with real scalars λ, 1-λ). So S ∈ D.

The map Φ is affine on D: ω\_S(a) = ⟨π\_φ(a)Sξ\_φ, ξ\_φ⟩ = ⟨π\_φ(a)(λT\_1 + (1-λ)T\_2)ξ\_φ, ξ\_φ⟩ = λ⟨π\_φ(a)T\_1ξ\_φ, ξ\_φ⟩ + (1-λ)⟨π\_φ(a)T\_2ξ\_φ, ξ\_φ⟩ = λω\_\{T\_1\}(a) + (1-λ)ω\_\{T\_2\}(a). This uses linearity of the inner product in the first argument (the scalars λ, 1-λ are real, so sesquilinearity reduces to linearity). This is a direct proof depending on Node 1.1.1 (well-definedness of Φ).

  Type: claim, Inference: assumption, Status: pending, Taint: unresolved

\item \textbf{Node 1.1.3:} (Image of Φ lies in F\_\{S(A)\}(φ) — amended). For T in the domain D, ω\_T belongs to the face F\_\{S(A)\}(φ).

Proof: Since 0 ≤ T ≤ ||T||·I (as T ∈ π\_φ(A)'\textasciicircum{}+ is a bounded positive operator), for any a ∈ A: ω\_T(a*a) = ⟨Tπ\_φ(a)ξ\_φ, π\_φ(a)ξ\_φ⟩ ≤ ||T||·⟨π\_φ(a)ξ\_φ, π\_φ(a)ξ\_φ⟩ = ||T||·φ(a*a). Hence ω\_T ≤ ||T||·φ as positive functionals.

A state ψ belongs to F\_\{S(A)\}(φ) iff there exists λ > 0 with ψ ≤ λφ. Here is the proof of this characterisation (both directions): (Forward) If ψ ∈ F\_\{S(A)\}(φ), then by definition of the generated face, φ = tψ + (1-t)σ for some σ ∈ S(A) and t ∈ (0,1). Since σ ≥ 0, we get tψ ≤ φ, i.e., ψ ≤ (1/t)φ. (Backward) If ψ ≤ λφ with λ ≥ 1 (λ ≥ 1 follows from ψ(1) = φ(1) = 1), define σ = (λφ - ψ)/(λ-1) when λ > 1. Then σ ≥ 0, σ(1) = (λ-1)/(λ-1) = 1, so σ ∈ S(A). And φ = (1/λ)ψ + (1-1/λ)σ, which writes φ as a proper convex combination with ψ, so ψ ∈ F\_\{S(A)\}(φ). When λ = 1, ψ = φ ∈ F\_\{S(A)\}(φ) trivially.

Since ω\_T ≤ ||T||·φ and ⟨Tξ\_φ, ξ\_φ⟩ = 1, we have ||T|| ≥ 1 > 0, so the characterisation gives ω\_T ∈ F\_\{S(A)\}(φ). This is a direct proof, not an assumption.

  Type: claim, Inference: assumption, Status: pending, Taint: unresolved

\item \textbf{Node 1.1.4:} (Injectivity of Φ). The map T ↦ ω\_T is injective. Proof: Suppose ω\_\{T\_1\} = ω\_\{T\_2\}. Then for all a ∈ A, ⟨π\_φ(a)(T\_1 - T\_2)ξ\_φ, ξ\_φ⟩ = 0. Since ξ\_φ is cyclic for π\_φ, the set \{π\_φ(a)ξ\_φ : a ∈ A\} is dense in H\_φ. Setting S = T\_1 - T\_2 ∈ π\_φ(A)' (self-adjoint), we have ⟨Sπ\_φ(a)ξ\_φ, ξ\_φ⟩ = 0 for all a (using S ∈ commutant to move S past π\_φ(a)). Then ⟨Sξ\_φ, π\_φ(a)*ξ\_φ⟩ = 0 for all a, i.e., Sξ\_φ ⊥ π\_φ(A)*ξ\_φ = π\_φ(A)ξ\_φ, which is dense. So Sξ\_φ = 0. Now for any b ∈ A: Sπ\_φ(b)ξ\_φ = π\_φ(b)Sξ\_φ = 0 (since S commutes with π\_φ). By density of π\_φ(A)ξ\_φ and continuity of S, S = 0, so T\_1 = T\_2.

  Type: claim, Inference: assumption, Status: validated, Taint: unresolved

\item \textbf{Node 1.1.5:} (Surjectivity of Φ via Radon-Nikodym — amended). The map Φ is surjective onto F\_\{S(A)\}(φ).

Proof: Let ψ ∈ F\_\{S(A)\}(φ). By the face characterisation (proved in Node 1.1.3): ψ ∈ F\_\{S(A)\}(φ) iff ψ ≤ λφ for some λ > 0. Specifically, φ = tψ + (1-t)σ for some t ∈ (0,1), giving ψ ≤ (1/t)φ with λ = 1/t.

By the Radon-Nikodym theorem for states (Alfsen-Shultz 2001, Proposition 5.31; Takesaki I.9): given ψ ≤ λφ, there exists a unique T ∈ π\_φ(A)' with 0 ≤ T ≤ λI such that ψ(a) = ⟨π\_φ(a)Tξ\_φ, ξ\_φ⟩ for all a ∈ A. This result is established for C*-algebras (where GNS gives bounded representations) and extends to *-algebras with Archimedean conditions (where GNS is also bounded, per Node 1.3.1).

Since ψ is a state: ψ(1) = 1 = ⟨Tξ\_φ, ξ\_φ⟩. Hence T is in the domain of Φ and Φ(T) = ψ.

This is a direct proof depending on Node 1.1.3 (face characterisation), the Radon-Nikodym theorem for states (external reference), and Node 1.3.1 (scope: GNS gives bounded representations).

  Type: claim, Inference: assumption, Status: pending, Taint: unresolved

\item \textbf{Node 1.1.6:} (Dimension consequence — amended). Since Φ is an affine isomorphism between F\_\{S(A)\}(φ) and D = \{T ∈ π\_φ(A)'\textasciicircum{}+ : ⟨Tξ\_φ, ξ\_φ⟩ = 1\}, we have dim(F\_\{S(A)\}(φ)) = dim(D). We claim dim(D) = dim\_R(π\_φ(A)'\textasciicircum{}\{sa\}) - 1.

Let V = π\_φ(A)'\textasciicircum{}\{sa\} and f : V → R the functional f(T) = ⟨Tξ\_φ, ξ\_φ⟩, which is real-linear. D = π\_φ(A)'\textasciicircum{}+ ∩ f⁻¹(1), a cross-section of the positive cone by a hyperplane.

The formula dim(D) = dim\_R(V) - 1 requires showing that aff(D) = f⁻¹(1), i.e., that the affine span of D fills the entire normalisation hyperplane. This follows from two facts:

(a) The cone π\_φ(A)'\textasciicircum{}+ SPANS V (every self-adjoint operator is a difference of two positive operators: T = T₊ - T₋ where T₊, T₋ ≥ 0).

(b) The cone π\_φ(A)'\textasciicircum{}+ has NONEMPTY INTERIOR in V: the identity I ∈ π\_φ(A)' is a strictly positive element, so a ball around I in V lies in π\_φ(A)'\textasciicircum{}+.

From (a) and (b): take T₀ = I/⟨ξ\_φ, ξ\_φ⟩ = I (since ξ\_φ is a unit vector), so f(T₀) = 1 and T₀ ∈ int(π\_φ(A)'\textasciicircum{}+). For any S ∈ V with f(S) = 0 (i.e., S ∈ ker(f)), T₀ + εS ∈ π\_φ(A)'\textasciicircum{}+ for sufficiently small ε > 0 (since T₀ is interior), and f(T₀ + εS) = 1. Hence T₀ + εS ∈ D, showing ker(f) ⊆ \{directions of D from T₀\}, so aff(D) = f⁻¹(1) and dim(D) = dim\_R(V) - 1.

Note: Faithfulness of f on π\_φ(A)'\textasciicircum{}+ (equivalently, ξ\_φ being separating for π\_φ(A)') is NOT required for this dimension formula — it suffices that f is a non-zero linear functional on V and the cone has interior. Faithfulness is used elsewhere (injectivity, Node 1.1.4).

Combining with Nodes 1.1.1–1.1.5: dim(F\_\{S(A)\}(φ)) = dim(D) = dim\_R(π\_φ(A)'\textasciicircum{}\{sa\}) - 1. This is a deduction depending on Nodes 1.1.1–1.1.5.

  Type: claim, Inference: assumption, Status: pending, Taint: unresolved

\end{enumerate}
\item \textbf{Node 1.2:} (Abstract Face-Dimension Bound — Weis-Shirokov). Let K be a compact convex set and let C = \{x ∈ K : f\_j(x) = β\_j, j = 1,...,m\} (or ≤ β\_j, or mixed) where each f\_j is a continuous affine functional (or more generally, a lower semicontinuous generalized affine map). For any x ∈ C: dim(F\_K(x)) ≤ dim(F\_C(x)) + m. In particular, if x is an extreme point of C (dim(F\_C(x)) = 0), then dim(F\_K(x)) ≤ m.

  Type: claim, Inference: assumption, Status: pending, Taint: unresolved

\begin{enumerate}
\item \textbf{Node 1.2.1:} (Single-constraint case — Weis-Shirokov Theorem 2, amended). Let K be a convex set in a locally convex space, f : K → R ∪ \{+∞\} a lower semicontinuous generalized affine map, and β ∈ R. Define the level set C\_= = \{x ∈ K : f(x) = β\} or sublevel set C\_≤ = \{x ∈ K : f(x) ≤ β\}. For any x ∈ C (= or ≤): dim(F\_K(x)) ≤ dim(F\_C(x)) + 1.

This is taken as an EXTERNAL RESULT: Weis \& Shirokov (2021), 'The face generated by a point, generalized affine constraints, and quantum theory', Theorem 2 (for level sets) and Theorem 4 (for sublevel sets). The proof is non-trivial and we cite it rather than reproduce it. The key steps in their argument are: (i) Restrict f to aff(F\_K(x)), where it remains affine. (ii) The level set \{f = f(x)\} ∩ aff(F\_K(x)) has codimension at most 1 in aff(F\_K(x)). (iii) The face F\_C(x) satisfies F\_C(x) ⊇ F\_K(x) ∩ C ∩ \{relative interior condition\}, giving dim(F\_C(x)) ≥ dim(F\_K(x)) - 1. (iv) The sublevel case requires additional argument via the exposed face structure and lower semicontinuity.

Note: The naive 'constraint propagation' argument (if y ∈ F\_K(x) then f(y) = β) does NOT work directly — for sublevel sets, f(y) may exceed β even when x = λy + (1-λ)z with f(x) ≤ β, because f(z) can compensate. Weis-Shirokov's proof handles this correctly through the face structure analysis. Compactness of K is not required for the statement; it is used in applications to guarantee existence of extreme points. Depends on external reference: Weis-Shirokov 2021.

  Type: claim, Inference: assumption, Status: pending, Taint: unresolved

\item \textbf{Node 1.2.2:} (Multi-constraint induction — Weis-Shirokov Corollary 5, amended). Under m constraints C = \{x ∈ K : f\_j(x) ≤ β\_j (or = β\_j), j = 1,...,m\}: dim(F\_K(x)) ≤ dim(F\_C(x)) + m. Proof: Apply Node 1.2.1 (single-constraint case) iteratively. Define K\_0 = K and for j = 1,...,m, define K\_j = \{x ∈ K\_\{j-1\} : f\_j(x) ≤ β\_j\} (for inequality constraints) or K\_j = \{x ∈ K\_\{j-1\} : f\_j(x) = β\_j\} (for equality constraints). Then C = K\_m. By Node 1.2.1 applied at each step: dim(F\_\{K\_\{j-1\}\}(x)) ≤ dim(F\_\{K\_j\}(x)) + 1. Summing (telescoping): dim(F\_K(x)) = dim(F\_\{K\_0\}(x)) ≤ dim(F\_\{K\_m\}(x)) + m = dim(F\_C(x)) + m. Note: each constraint costs at most 1 dimension whether it is an equality or inequality, since Node 1.2.1 handles both level and sublevel sets. This is a deduction by iteration from Node 1.2.1.

  Type: claim, Inference: assumption, Status: pending, Taint: unresolved

\item \textbf{Node 1.2.3:} (Extremal specialisation — amended). If x is an extreme point of C (i.e., F\_C(x) = \{x\}, so dim(F\_C(x)) = 0), the multi-constraint bound (Node 1.2.2) gives dim(F\_K(x)) ≤ 0 + m = m. This is the key input for the main theorem: taking K = S(A) and C the affinely constrained state space with m constraints (with hypotheses verified in Node 1.3.1, confirming that the constraints f\_j(ψ) = ψ(a\_j) are continuous affine on S(A)) yields dim(F\_\{S(A)\}(φ)) ≤ m for extreme points φ of C. This is a deduction by universal instantiation from Node 1.2.2, depending on Nodes 1.2.2 and 1.2.4.

  Type: claim, Inference: assumption, Status: pending, Taint: unresolved

\item \textbf{Node 1.2.4:} (Generalized vs ordinary affine maps — amended). A continuous affine functional f : K → R (where K is a compact convex set in a locally convex space) is a fortiori lower semicontinuous, and since dom(f) = K (no +∞ values), it is a generalized affine map in the sense of Weis-Shirokov. Hence for bounded self-adjoint elements a\_j ∈ A\textasciicircum{}\{sa\} (where A is a C*-algebra or Archimedean order unit space), the maps f\_j : S(A) → R defined by f\_j(φ) = φ(a\_j) are continuous affine functionals — continuous by definition of the weak-* topology (evaluation at a fixed element is weak-* continuous by definition) — and therefore satisfy the hypotheses of Node 1.2.

The generalized version (allowing f : K → R ∪ \{+∞\} with f lower semicontinuous) is needed for unbounded observables in infinite-dimensional quantum mechanics: for an unbounded self-adjoint operator H affiliated with B(H), f(ρ) = Tr(Hρ) may equal +∞ when ρ is a density operator not in the form domain of H. This generalization is used in Node 1.8.3 (energy constraints). This is a deduction depending on the definitions of generalized\_affine\_map, state\_space, and self\_adjoint\_part.

  Type: claim, Inference: assumption, Status: pending, Taint: unresolved

\end{enumerate}
\item \textbf{Node 1.3:} (Main Dimension Bound). Let A be a unital *-algebra, a\_1,...,a\_m ∈ A\textasciicircum{}\{sa\}, β\_1,...,β\_m ∈ R, and C = \{φ ∈ S(A) : φ(a\_j) = β\_j\}. If φ is an extreme point of C, then dim\_R(π\_φ(A)'\textasciicircum{}\{sa\}) ≤ m + 1. Proof: By Node 1.1 (face-commutant correspondence), dim(F\_\{S(A)\}(φ)) = dim\_R(π\_φ(A)'\textasciicircum{}\{sa\}) - 1. By Node 1.2 (Weis-Shirokov), since φ is extreme in C, dim(F\_\{S(A)\}(φ)) ≤ m. Combining: dim\_R(π\_φ(A)'\textasciicircum{}\{sa\}) - 1 ≤ m, hence dim\_R(π\_φ(A)'\textasciicircum{}\{sa\}) ≤ m + 1.

  Type: claim, Inference: assumption, Status: pending, Taint: unresolved

\begin{enumerate}
\item \textbf{Node 1.3.1:} (Hypotheses verification — CORRECTED scope, 2nd amendment). The main theorem applies to: (A) C*-algebras, (B) Archimedean order unit spaces (V, V\textasciicircum{}+, e) with order-unit norm, and (C) *-algebras equipped with an Archimedean quadratic module M. For each case, we verify the hypotheses of Nodes 1.1 (face-commutant) and 1.2 (Weis-Shirokov). This node depends on Nodes 1.1 and 1.2.

(i) GNS gives BOUNDED representations: For C*-algebras, the C*-identity ||a*a|| = ||a||² gives ||π(a)|| ≤ ||a||. For Archimedean order unit spaces, the order-unit norm provides boundedness. For *-algebras with Archimedean quadratic module M, M-positive states are uniformly bounded (Cauchy-Schwarz + Archimedean: |φ(a)| ≤ √N\_a where a*a ≤ N\_a·1 mod M). In all three cases, the GNS representation is bounded, so the commutant π\_φ(A)' = \{T ∈ B(H) : Tπ(a) = π(a)T ∀a\} is a well-defined von Neumann algebra, and the face-commutant correspondence (Node 1.1) applies.

(ii) Compactness of K = S(A): For C*-algebras, states lie in the unit ball of A* which is weak-* compact by Banach-Alaoglu. For Archimedean order unit spaces, the order-unit norm gives a Banach space structure and Banach-Alaoglu applies. For *-algebras with Archimedean quadratic module, S\_M is compact in the pointwise topology by Tychonoff's theorem (S\_M embeds as a closed subset of a product of compact intervals, using the Archimedean bounds). In all cases K = S(A) (or S\_M) is compact convex.

(iii) Continuity of constraint functionals: The functionals f\_j(ψ) = ψ(a\_j) are continuous affine on S(A) by DEFINITION of the weak-* (or pointwise) topology — evaluation at a fixed element is continuous by the very definition of this topology. This holds regardless of whether A has a norm. The Weis-Shirokov theorem (Node 1.2) applies with K = S(A) and m constraints.

Note: For GENERAL *-algebras without the Archimedean property, the GNS construction may give unbounded operators, the commutant is not well-defined in the standard sense, and the state space need not be compact. The theorem does not apply in that generality. This node is a verification/deduction step, not an assumption.

  Type: claim, Inference: assumption, Status: pending, Taint: unresolved

\item \textbf{Node 1.3.2:} (Step 1: Face-commutant gives geometric-algebraic bridge — amended). By Node 1.3.1, the hypotheses of Node 1.1 are satisfied (A is in the applicable scope, GNS yields bounded representations). Therefore, by Node 1.1 (Alfsen-Shultz face-commutant correspondence): dim(F\_\{S(A)\}(φ)) = dim\_R(π\_φ(A)'\textasciicircum{}\{sa\}) - 1. This converts the geometric quantity (face dimension in the state space) into the algebraic quantity (real dimension of the self-adjoint part of the commutant). This is a deduction by universal instantiation of Node 1.1, depending on Nodes 1.1 and 1.3.1.

  Type: claim, Inference: assumption, Status: pending, Taint: unresolved

\item \textbf{Node 1.3.3:} (Step 2: Extremality implies face bound — amended). Since φ is an extreme point of C, by definition F\_C(φ) = \{φ\}, hence dim(F\_C(φ)) = 0. By Node 1.3.1, the hypotheses of Node 1.2 (Weis-Shirokov) are satisfied with K = S(A) compact convex and f\_j(ψ) = ψ(a\_j) continuous affine. Applying Node 1.2.3 (extremal specialisation of Weis-Shirokov) with m equality constraints f\_j(ψ) = β\_j: dim(F\_\{S(A)\}(φ)) ≤ dim(F\_C(φ)) + m = 0 + m = m. The same argument applies to inequality constraints φ(a\_j) ≤ β\_j or mixed constraints, since Node 1.2 covers sublevel sets C = \{x ∈ K : f(x) ≤ β\} as well as level sets (see Node 1.3.5 for the mixed case). This is a deduction depending on Nodes 1.2 (specifically 1.2.3), 1.3.1, and the definition of extreme point.

  Type: claim, Inference: assumption, Status: pending, Taint: unresolved

\item \textbf{Node 1.3.4:} (Step 3: Algebraic conclusion). Combining Steps 1 and 2: dim\_R(π\_φ(A)'\textasciicircum{}\{sa\}) - 1 = dim(F\_\{S(A)\}(φ)) ≤ m. Therefore dim\_R(π\_φ(A)'\textasciicircum{}\{sa\}) ≤ m + 1. This completes the proof of part (I) of the main theorem.

  Type: claim, Inference: assumption, Status: validated, Taint: unresolved

\item \textbf{Node 1.3.5:} (Extension to inequality and mixed constraints). The same argument works for inequality constraints φ(a\_j) ≤ β\_j or mixed equality/inequality: the Weis-Shirokov theorem (Node 1.2) applies to both sublevel sets and level sets, and the inductive argument treats each constraint independently, whether it is an equality or inequality. Hence the bound dim\_R(π\_φ(A)'\textasciicircum{}\{sa\}) ≤ m + 1 holds for all three constraint types.

  Type: claim, Inference: assumption, Status: validated, Taint: unresolved

\end{enumerate}
\item \textbf{Node 1.4:} (Commutant Decomposition and Multiplicity Bound). Under the hypotheses of the Main Dimension Bound (Node 1.3), if dim\_R(π\_φ(A)'\textasciicircum{}\{sa\}) < ∞, then the commutant π\_φ(A)' is a finite-dimensional von Neumann algebra, hence by Artin-Wedderburn π\_φ(A)' ≅ ⊕\_\{j=1\}\textasciicircum{}k M\_\{n\_j\}(C) for pairwise inequivalent irreducible summands σ\_j with multiplicities n\_j. Therefore dim\_R(π\_φ(A)'\textasciicircum{}\{sa\}) = Σ\_\{j=1\}\textasciicircum{}k n\_j², and the bound becomes Σ\_\{j=1\}\textasciicircum{}k n\_j² ≤ m + 1.

  Type: claim, Inference: assumption, Status: pending, Taint: unresolved

\begin{enumerate}
\item \textbf{Node 1.4.1:} (Finite-dimensionality of commutant). From Node 1.3, if φ is extreme in C under m constraints, then dim\_R(π\_φ(A)'\textasciicircum{}\{sa\}) ≤ m + 1 < ∞. Since π\_φ(A)'\textasciicircum{}\{sa\} is the self-adjoint part of the *-algebra π\_φ(A)', and every element of π\_φ(A)' decomposes as T = Re(T) + i·Im(T) with Re(T), Im(T) ∈ π\_φ(A)'\textasciicircum{}\{sa\}, we have dim\_C(π\_φ(A)') = dim\_R(π\_φ(A)'\textasciicircum{}\{sa\}) ≤ m + 1 < ∞. Thus π\_φ(A)' is a finite-dimensional von Neumann algebra (it is a *-subalgebra of B(H\_φ) closed in all operator topologies, being finite-dimensional).

  Type: claim, Inference: assumption, Status: validated, Taint: unresolved

\item \textbf{Node 1.4.2:} (Artin-Wedderburn decomposition of commutant — amended). By Node 1.4.1, π\_φ(A)' is a finite-dimensional C*-algebra (finite-dimensional von Neumann algebra). By the Artin-Wedderburn theorem (Lam, 'A First Course in Noncommutative Rings', Theorem 3.5; or Kadison-Ringrose Vol. I, Theorem 6.6.5 for the C*-version), combined with the Gelfand-Mazur theorem (the only finite-dimensional complex division algebra is C): π\_φ(A)' ≅ M\_\{n\_1\}(C) ⊕ M\_\{n\_2\}(C) ⊕ ... ⊕ M\_\{n\_k\}(C) for unique integers n\_1, ..., n\_k ≥ 1 (up to reordering). Here k is the number of minimal central projections in π\_φ(A)'. This is an application of the Artin-Wedderburn theorem (external reference), depending on Node 1.4.1.

  Type: claim, Inference: assumption, Status: pending, Taint: unresolved

\item \textbf{Node 1.4.3:} (Correspondence with representation decomposition — amended). The Artin-Wedderburn blocks of π\_φ(A)' (from Node 1.4.2) correspond to the isotypic components of π\_φ. Since π\_φ(A)' is finite-dimensional (Node 1.4.1), the representation π\_φ is completely reducible: the minimal projections in each block M\_\{n\_j\}(C) of the commutant project onto irreducible subrepresentations, and their orthogonal complements are invariant. This gives π\_φ ≅ ⊕\_\{j=1\}\textasciicircum{}k σ\_j\textasciicircum{}\{⊕n\_j\}, where σ\_1, ..., σ\_k are pairwise unitarily inequivalent irreducible representations.

The proof uses the commutant theorem (bicommutant + Schur's lemma): (1) The minimal central projections P\_j in π\_φ(A)' decompose H\_φ = ⊕ P\_j(H\_φ). (2) On each P\_j(H\_φ), π\_φ(A)' restricts to the factor M\_\{n\_j\}(C), so by the bicommutant theorem π\_φ(A)'' restricts to B(H\_\{σ\_j\}) ⊗ Id\_\{n\_j\}. (3) Complete reducibility from (above) decomposes P\_j(H\_φ) ≅ H\_\{σ\_j\} ⊗ C\textasciicircum{}\{n\_j\} with π\_φ = σ\_j ⊗ Id and π\_φ(A)' = Id ⊗ M\_\{n\_j\}(C). The j-th block M\_\{n\_j\}(C) acts on the multiplicity space C\textasciicircum{}\{n\_j\} by taking linear combinations of the n\_j copies: (T·(v\_1,...,v\_\{n\_j\}))\_p = Σ\_q t\_\{pq\} v\_q. This is a theorem (not an assumption) depending on Nodes 1.4.1 and 1.4.2.

  Type: claim, Inference: assumption, Status: pending, Taint: unresolved

\item \textbf{Node 1.4.4:} (Dimension calculation — amended). By Node 1.4.2 (Artin-Wedderburn): π\_φ(A)' ≅ ⊕\_\{j=1\}\textasciicircum{}k M\_\{n\_j\}(C). The self-adjoint part of each block M\_\{n\_j\}(C)\textasciicircum{}\{sa\} consists of n\_j × n\_j Hermitian matrices, a real vector space of dimension n\_j². Therefore: dim\_R(π\_φ(A)'\textasciicircum{}\{sa\}) = Σ\_\{j=1\}\textasciicircum{}k dim\_R(M\_\{n\_j\}(C)\textasciicircum{}\{sa\}) = Σ\_\{j=1\}\textasciicircum{}k n\_j². Combining with Node 1.3 (dim\_R(π\_φ(A)'\textasciicircum{}\{sa\}) ≤ m + 1): Σ\_\{j=1\}\textasciicircum{}k n\_j² ≤ m + 1. Combined with Node 1.4.3 (the n\_j are the representation multiplicities), this yields part (II) of the main theorem. This is a computation/deduction depending on Nodes 1.4.2, 1.4.3, and 1.3.

  Type: claim, Inference: assumption, Status: pending, Taint: unresolved

\item \textbf{Node 1.4.5:} (Contrapositive for infinite-dimensional case — amended). If π\_φ(A)' is infinite-dimensional (dim\_R(π\_φ(A)'\textasciicircum{}\{sa\}) = ∞), then φ cannot be an extreme point of any set C defined by finitely many CONTINUOUS affine constraints. This is the contrapositive of Node 1.3: if dim\_R(π\_φ(A)'\textasciicircum{}\{sa\}) > m + 1 for any finite m, then φ is not extreme in C. Equivalently (by Nodes 1.4.1–1.4.4), under finitely many continuous affine constraints, every extreme point has a finite-dimensional commutant and the multiplicity bound Σ n\_j² ≤ m + 1 applies. This is the useful form for infinite-dimensional algebras such as B(H) (Node 1.8) where infinite-dimensional commutants exist. This is a deduction by contrapositive from Node 1.3, depending on Nodes 1.3, 1.4.1–1.4.4.

  Type: claim, Inference: assumption, Status: pending, Taint: unresolved

\end{enumerate}
\item \textbf{Node 1.5:} (Recovery: Abelian Case → Carathéodory). When A = C(X) for a compact Hausdorff space X, the main theorem (Σ n\_j² ≤ m+1) recovers Carathéodory's theorem for moment problems: every extreme point of the set of probability measures on X satisfying m moment constraints is a discrete measure supported on at most m+1 points.

  Type: claim, Inference: assumption, Status: pending, Taint: unresolved

\begin{enumerate}
\item \textbf{Node 1.5.1:} (Abelian algebra setup). Let A = C(X) for a compact Hausdorff space X. The involution is pointwise conjugation f*(x) = f(x)̄, and A\textasciicircum{}\{sa\} = C(X, R) (real-valued continuous functions). By Riesz representation, S(C(X)) is the set of Borel probability measures on X: φ ↦ μ\_φ where φ(f) = ∫\_X f dμ\_φ.

  Type: claim, Inference: assumption, Status: validated, Taint: unresolved

\item \textbf{Node 1.5.2:} (All irreducible representations are 1-dimensional and all GNS multiplicities are 1). For A = C(X) with X compact Hausdorff, we establish two claims: (A) Every irreducible representation of C(X) is 1-dimensional. (B) For any state phi on C(X), all multiplicities in the GNS decomposition are n\_j = 1.

Proof of (A): Let pi : C(X) -> B(H) be an irreducible *-representation. Since C(X) is abelian, pi(C(X)) is an abelian *-subalgebra of B(H). By Schur's lemma (Takesaki I, Proposition I.9.1; or equivalently, for an irreducible representation, the commutant pi(A)' = C*I), every operator in pi(C(X)) commutes with every operator in B(H) that commutes with pi(C(X)), and since pi(C(X)) is abelian, pi(C(X)) is contained in its own commutant pi(C(X))' = C*I. Hence every operator in pi(C(X)) is a scalar multiple of I, and irreducibility forces dim(H) = 1. Therefore pi = ev\_x for some x in X (a character/point evaluation). This identification of irreducible representations of C(X) with point evaluations is the content of the Gelfand-Naimark theorem for commutative C*-algebras (Kadison-Ringrose I, Theorem 4.4.3): the maximal ideal space of C(X) is homeomorphic to X.

Proof of (B): For a state phi on C(X), the GNS representation pi\_phi is cyclic by construction. Since C(X) is abelian, pi\_phi(C(X)) is abelian, hence pi\_phi(C(X)) is contained in pi\_phi(C(X))' (the commutant). It follows that pi\_phi(C(X))' is also abelian: indeed, if the commutant decomposes (via Artin-Wedderburn, once known to be finite-dimensional) as pi\_phi(C(X))' = direct\_sum M\_\{n\_j\}(C), then pi\_phi(C(X)) is contained in this commutant, and pi\_phi(C(X)) contains the identity of each block (by irreducible decomposition). But pi\_phi(C(X)) being abelian and contained in the commutant forces each M\_\{n\_j\}(C) to be abelian, hence n\_j = 1 for all j. Alternatively, without assuming finite-dimensionality: for a cyclic representation of a commutative C*-algebra, the commutant pi\_phi(C(X))' equals L\textasciicircum{}infty(X, mu\_phi) acting by multiplication (a maximal abelian von Neumann algebra), which is abelian. In either case, all multiplicities satisfy n\_j = 1.

This is a deduction depending on Node 1.5.1 (abelian algebra setup), Schur's lemma (external: Takesaki I, Prop. I.9.1), and the Gelfand-Naimark theorem for commutative C*-algebras (external: Kadison-Ringrose I, Thm 4.4.3).

  Type: claim, Inference: assumption, Status: pending, Taint: unresolved

\item \textbf{Node 1.5.3:} (GNS decomposition for states on C(X): finite-dimensional commutant implies finite atomic support, and explicit GNS computation). Let A = C(X) for X compact Hausdorff. For a state phi\_mu on C(X) corresponding to a Borel probability measure mu on X (via Node 1.5.1), we establish two results:

(I) FINITE COMMUTANT DIMENSION IMPLIES FINITE ATOMIC SUPPORT: If dim\_R(pi\_mu(C(X))'\textasciicircum{}\{sa\}) < infinity, then mu is atomic with finite support. Proof: The GNS commutant pi\_mu(C(X))' is isomorphic (as a von Neumann algebra) to L\textasciicircum{}infty(X, mu) acting by multiplication on L\textasciicircum{}2(X, mu). The algebra L\textasciicircum{}infty(X, mu) is finite-dimensional if and only if mu is a finite sum of point masses (atoms). Forward direction: if L\textasciicircum{}infty(X, mu) is finite-dimensional, say of dimension k, then L\textasciicircum{}2(X, mu) is k-dimensional, which forces mu to be supported on exactly k atoms (the indicator functions of the atoms form a basis). Backward direction: if mu = sum\_\{j=1\}\textasciicircum{}k lambda\_j delta\_\{x\_j\} with distinct x\_j and lambda\_j > 0, then L\textasciicircum{}infty(X, mu) = C\textasciicircum{}k (functions on a k-point set). Thus: dim\_R(pi\_mu(C(X))'\textasciicircum{}\{sa\}) < infinity implies mu = sum\_\{j=1\}\textasciicircum{}k lambda\_j delta\_\{x\_j\} for some k < infinity with distinct points x\_1, ..., x\_k in X and weights lambda\_j > 0 summing to 1.

(II) EXPLICIT GNS DECOMPOSITION FOR FINITELY-SUPPORTED MEASURES: For mu = sum\_\{j=1\}\textasciicircum{}k lambda\_j delta\_\{x\_j\} with distinct points x\_1, ..., x\_k in X and lambda\_j > 0, the GNS representation is pi\_mu = direct\_sum\_\{j=1\}\textasciicircum{}k ev\_\{x\_j\} (each evaluation representation appearing with multiplicity 1), and the commutant is pi\_mu(C(X))' = C\textasciicircum{}k. Hence dim\_R(pi\_mu(C(X))'\textasciicircum{}\{sa\}) = k (each summand contributes n\_j\textasciicircum{}2 = 1\textasciicircum{}2 = 1, where n\_j = 1 is the multiplicity of the j-th evaluation representation).

Proof of (II): The GNS Hilbert space is H\_mu = completion of C(X)/N\_mu where N\_mu = \{f : phi\_mu(f*f) = 0\} = \{f : sum lambda\_j |f(x\_j)|\textasciicircum{}2 = 0\} = \{f : f(x\_j) = 0 for all j = 1,...,k\}. Thus C(X)/N\_mu is isomorphic to C\textasciicircum{}k via f -> (f(x\_1), ..., f(x\_k)), and H\_mu = C\textasciicircum{}k (already complete since finite-dimensional). The inner product is the weighted one: <f, g>\_mu = sum lambda\_j f(x\_j) conj(g(x\_j)). The GNS representation acts by multiplication: pi\_mu(f) = diag(f(x\_1), ..., f(x\_k)). By Urysohn's lemma (X compact Hausdorff implies X normal, and \{x\_1,...,x\_k\} is finite hence closed), for any values (alpha\_1,...,alpha\_k) in C\textasciicircum{}k there exists g in C(X) with g(x\_j) = alpha\_j, so pi\_mu(C(X)) is the FULL diagonal subalgebra of M\_k(C). The unitary U : (C\textasciicircum{}k, weighted) -> (C\textasciicircum{}k, standard) defined by U(e\_j) = sqrt(lambda\_j) e\_j intertwines pi\_mu with direct\_sum\_\{j=1\}\textasciicircum{}k ev\_\{x\_j\}, verifying the decomposition. The commutant of the full diagonal subalgebra in M\_k(C) is the diagonal subalgebra itself, so pi\_mu(C(X))' = C\textasciicircum{}k (abelian, confirming n\_j = 1 for all j by Node 1.5.2).

This is a deduction depending on Nodes 1.5.1 (states = probability measures), 1.5.2 (all multiplicities n\_j = 1 for abelian case), the GNS construction, and Urysohn's lemma (X compact Hausdorff implies normal, hence continuous functions separate finite point sets).

  Type: claim, Inference: assumption, Status: pending, Taint: unresolved

\item \textbf{Node 1.5.4:} (Applying the main bound: recovery of the Richter-Rogosinski theorem). Let A = C(X) with X compact Hausdorff, and let phi be a state on C(X) corresponding to a probability measure mu on X (Node 1.5.1). Suppose phi is an extreme point of the constrained set C = \{psi in S(C(X)) : psi(f\_j) = beta\_j, j = 1,...,m\} for f\_j in C(X, R).

Step 1 (Finite-dimensionality of commutant): By the main bound (Node 1.3), dim\_R(pi\_phi(C(X))'\textasciicircum{}\{sa\}) <= m + 1 < infinity. Hence the commutant is finite-dimensional.

Step 2 (Finite support): By Node 1.5.3, part (I), the finite-dimensionality of pi\_mu(C(X))' = L\textasciicircum{}infty(X, mu) forces mu to be atomic with finite support: mu = sum\_\{j=1\}\textasciicircum{}k lambda\_j delta\_\{x\_j\} for some k < infinity, with distinct points x\_1,...,x\_k in X and weights lambda\_j > 0.

Step 3 (GNS decomposition and multiplicity bound): By Node 1.5.3, part (II), the GNS representation is pi\_mu = direct\_sum\_\{j=1\}\textasciicircum{}k ev\_\{x\_j\} with each multiplicity n\_j = 1 (as established in Node 1.5.2). By Node 1.4 (commutant decomposition), dim\_R(pi\_mu(C(X))'\textasciicircum{}\{sa\}) = sum\_\{j=1\}\textasciicircum{}k n\_j\textasciicircum{}2 = k * 1 = k. Combining with the main bound: k <= m + 1.

Conclusion: Every extreme point of the constrained set of probability measures \{mu in M\_1\textasciicircum{}+(X) : integral f\_j d mu = beta\_j, j = 1,...,m\} is a discrete (atomic) measure supported on at most m + 1 points. This is the Richter-Rogosinski theorem (Richter 1957, Rogosinski 1958), also known as the measure-theoretic form of Caratheodory's theorem for moment problems. Note that the logical flow is: main bound => finite commutant => finite support (via L\textasciicircum{}infty characterisation) => k <= m+1. The finite support is CONCLUDED (not assumed), resolving the circularity concern.

This is an application of the main theorem (Nodes 1.3, 1.4) to the abelian specialisation (Nodes 1.5.1, 1.5.2, 1.5.3). It is a deduction, not an assumption. Dependencies: Nodes 1.3, 1.4, 1.5.1, 1.5.2, 1.5.3.

  Type: claim, Inference: assumption, Status: pending, Taint: unresolved

\end{enumerate}
\item \textbf{Node 1.6:} (Recovery: Matrix Algebra → Classical Barvinok-Pataki). When A = M\_n(C), the main theorem recovers the classical Barvinok-Pataki bound for complex SDP: an extreme point of \{ρ ∈ S(M\_n(C)) : Tr(a\_j ρ) = β\_j\} has rank r satisfying r² ≤ m+1. The relationship to the standard formulation r² ≤ m (for the Hermitian PSD cone) is explained by the trace normalisation: the state space S(M\_n(C)) already incorporates Tr(ρ)=1, which in the standard SDP formulation is counted among the m constraints.

  Type: claim, Inference: assumption, Status: pending, Taint: unresolved

\begin{enumerate}
\item \textbf{Node 1.6.1:} (Matrix algebra setup). Let A = M\_n(C), the algebra of n×n complex matrices with involution a* = a† (conjugate transpose). Then A\textasciicircum{}\{sa\} = H\_n (the n² dimensional real space of Hermitian matrices). The state space S(M\_n(C)) = \{ρ ∈ H\_n : ρ ≥ 0, Tr(ρ) = 1\} is the set of n×n density matrices. States act by φ\_ρ(a) = Tr(ρa).

  Type: claim, Inference: assumption, Status: validated, Taint: unresolved

\item \textbf{Node 1.6.2:} (Unique irreducible representation). M\_n(C) is a simple algebra: it has exactly one irreducible representation (up to unitary equivalence), namely the identity representation id : M\_n(C) → B(C\textasciicircum{}n). Therefore every *-representation decomposes as π ≅ id\textasciicircum{}\{⊕r\} for some multiplicity r.

  Type: claim, Inference: assumption, Status: validated, Taint: unresolved

\item \textbf{Node 1.6.3:} (GNS of a rank-r density matrix). For a density matrix ρ of rank r, the GNS representation π\_ρ is unitarily equivalent to id\textasciicircum{}\{⊕r\} acting on C\textasciicircum{}n ⊗ C\textasciicircum{}r (or equivalently on ⊕\textasciicircum{}r C\textasciicircum{}n). The commutant is π\_ρ(M\_n(C))' ≅ M\_r(C) (the r×r matrices acting on the multiplicity space C\textasciicircum{}r). The self-adjoint part has dim\_R(M\_r(C)\textasciicircum{}\{sa\}) = r² (the space of r×r Hermitian matrices).

  Type: claim, Inference: assumption, Status: validated, Taint: unresolved

\item \textbf{Node 1.6.4:} (Applying the main bound). Let C = \{rho in S(M\_n(C)) : Tr(a\_j rho) = beta\_j, j=1,...,m\} where a\_j in H\_n are self-adjoint. Suppose rho is an EXTREME POINT of C. By Node 1.6.1, states on M\_n(C) are density matrices and phi\_rho(a) = Tr(rho a). By Node 1.6.2, M\_n(C) has exactly one irreducible representation (k=1). By Node 1.6.3, for rank(rho) = r the GNS representation gives pi\_rho = id\textasciicircum{}\{+r\}, so n\_1 = r, and the commutant satisfies dim\_R(pi\_rho(M\_n(C))'\textasciicircum{}\{sa\}) = r\textasciicircum{}2. By Node 1.3 (Main Dimension Bound) applied to this extreme point, r\textasciicircum{}2 = Sigma\_\{j=1\}\textasciicircum{}k n\_j\textasciicircum{}2 <= m+1 (with k=1, n\_1=r). Hence rank(rho) = r <= sqrt(m+1). [Inference: modus\_ponens. Dependencies: 1.3, 1.4, 1.6.1, 1.6.2, 1.6.3.]

  Type: claim, Inference: assumption, Status: pending, Taint: unresolved

\item \textbf{Node 1.6.5:} (Reconciliation with classical Barvinok-Pataki: the +1 offset). The classical BP for complex SDP states r² ≤ m for \{X ∈ H\_n\textasciicircum{}+ : Tr(A\_j X) = b\_j, j=1,...,m\}. Our bound gives r² ≤ m+1. The discrepancy is exactly 1, explained as follows: In the classical formulation, the trace constraint Tr(X) = 1 (or Tr(ρ) = 1) is typically counted as one of the m constraints. In our formulation, the state space S(M\_n(C)) already incorporates Tr(ρ) = 1, and the m constraints are additional. Thus our 'm' equals the classical 'm minus 1' (or equivalently, the classical m includes our m plus the trace constraint). Setting m\_classical = m\_ours + 1 gives r² ≤ m\_ours + 1 = m\_classical, recovering the standard bound.

  Type: claim, Inference: assumption, Status: validated, Taint: unresolved

\item \textbf{Node 1.6.6:} (Real vs complex subtlety — CORRECTED). For the real PSD cone S\textasciicircum{}n\_+ (real symmetric matrices), the classical BP bound is r(r+1)/2 ≤ m, strictly better than r² ≤ m. This arises because faces of S\textasciicircum{}n\_+ have dimension r(r+1)/2 (the dimension of the space of r×r real symmetric matrices), while faces of H\_n\textasciicircum{}+ have dimension r² (the dimension of r×r Hermitian matrices). Our theorem, working with complex *-algebras, naturally gives the Hermitian/complex version.

For the real case, one works with the real *-algebra A = M\_n(R) equipped with the TRANSPOSE involution a* = aᵀ (NOT the identity — the transpose is a genuine anti-automorphism). The self-adjoint part A\textasciicircum{}\{sa\} = S\_n (real symmetric matrices) is the order unit space one optimizes over, but M\_n(R) itself is the algebra. The GNS construction for real *-algebras (cf. Goodearl, 'Notes on real and complex C*-algebras'; Li, 'Real operator algebras') produces representations on REAL Hilbert spaces. For A = M\_n(R) with a rank-r state (via a rank-r PSD symmetric matrix), the GNS representation gives π ≅ id\textasciicircum{}\{⊕r\} on R\textasciicircum{}n ⊗ R\textasciicircum{}r, with commutant π(M\_n(R))' ≅ M\_r(R). The self-adjoint part under transpose is the space of r×r real symmetric matrices, giving dim\_R(π(A)'\textasciicircum{}\{sa\}) = r(r+1)/2, recovering the tighter real bound.

The full classification by real division algebras (Hurwitz theorem) yields THREE cases: (1) Real: A = M\_n(R), involution = transpose, self-adjoint part has dim r(r+1)/2, giving r(r+1)/2 ≤ m+1; (2) Complex: A = M\_n(C), involution = conjugate transpose, dim r², giving r² ≤ m+1; (3) Quaternionic: A = M\_n(H), involution = conjugate transpose, self-adjoint (self-dual) part has dim r(2r-1), giving r(2r-1) ≤ m+1. All three cases appear in Barvinok's original paper (2001). Our operator-algebraic framework via order unit spaces is well-suited to treat all three cases uniformly through the appropriate choice of *-algebra and involution.

  Type: claim, Inference: assumption, Status: pending, Taint: unresolved

\end{enumerate}
\item \textbf{Node 1.7:} (Recovery: Block-Diagonal Algebras → Mixed Bound). When A = ⊕\_\{i=1\}\textasciicircum{}N M\_\{d\_i\}(C), the main theorem gives the new bound Σ\_\{i: λ\_i > 0\} r\_i² ≤ m+1, where r\_i = rank(ρ\_i) for the density matrix in each active block. This simultaneously constrains ranks within blocks and the number of active blocks, interpolating between Carathéodory (all d\_i = 1) and Barvinok-Pataki (N = 1).

  Type: claim, Inference: assumption, Status: pending, Taint: unresolved

\begin{enumerate}
\item \textbf{Node 1.7.1:} (Block-diagonal algebra setup). Let A = ⊕\_\{i=1\}\textasciicircum{}N M\_\{d\_i\}(C). A state φ on A corresponds to a tuple (λ\_1 ρ\_1, ..., λ\_N ρ\_N) where λ\_i ≥ 0 with Σ\_\{i=1\}\textasciicircum{}N λ\_i = 1 are classical weights and ρ\_i is a density matrix on M\_\{d\_i\}(C) for each block with λ\_i > 0. The constraint φ(1) = 1 gives Σ λ\_i = 1. A self-adjoint element a ∈ A\textasciicircum{}\{sa\} is a tuple (a\_1, ..., a\_N) with a\_i ∈ M\_\{d\_i\}(C)\textasciicircum{}\{sa\}, and φ(a) = Σ λ\_i Tr(ρ\_i a\_i).

  Type: claim, Inference: assumption, Status: validated, Taint: unresolved

\item \textbf{Node 1.7.2:} (GNS decomposition for block-diagonal algebra). Let A = direct\_sum\_\{i=1\}\textasciicircum{}N M\_\{d\_i\}(C) with notation as in Node 1.7.1. For a state phi with active blocks I = \{i : lambda\_i > 0\}, the GNS representation decomposes as pi\_phi = direct\_sum\_\{i in I\} id\_\{d\_i\}\textasciicircum{}\{+r\_i\} where r\_i = rank(rho\_i). Proof sketch: By the GNS construction applied block-by-block, the GNS of the restriction of phi to the i-th summand M\_\{d\_i\}(C) (with weight lambda\_i > 0) yields id\_\{d\_i\}\textasciicircum{}\{+r\_i\} on C\textasciicircum{}\{d\_i\} tensor C\textasciicircum{}\{r\_i\} (as in Node 1.6.3 applied to each block). The GNS of a direct sum state is the direct sum of the GNS representations of each active block. The irreducible representations of A = direct\_sum M\_\{d\_i\}(C) are exactly pi\_i := id\_\{d\_i\} composed with the canonical projection pr\_i : A -> M\_\{d\_i\}(C). These are pairwise inequivalent for distinct i, as can be seen by considering the central projection e\_i = (0,...,0,I\_\{d\_i\},0,...,0) in A: pi\_i(e\_i) = I\_\{d\_i\} (nonzero) while pi\_j(e\_i) = 0 for j \textbackslash\{\}!= i, so ker(pi\_i) \textbackslash\{\}!= ker(pi\_j), which precludes unitary equivalence. By Schur's lemma applied to the direct sum of pairwise inequivalent irreps with multiplicities r\_i, the commutant is pi\_phi(A)' = direct\_sum\_\{i in I\} M\_\{r\_i\}(C). [Inference: modus\_ponens. Dependencies: 1.7.1, 1.6.3. External: Wedderburn-Artin theory, Schur's lemma, GNS construction.]

  Type: claim, Inference: assumption, Status: pending, Taint: unresolved

\item \textbf{Node 1.7.3:} (Multiplicity bound for block-diagonal case). dim\_R(π\_φ(A)'\textasciicircum{}\{sa\}) = Σ\_\{i ∈ I\} r\_i². The main theorem gives Σ\_\{i ∈ I\} r\_i² ≤ m + 1. This is a simultaneous constraint: it bounds both max\_i r\_i (the maximum rank in any single block) and |I| (the number of active blocks). For example, if all r\_i = 1 then |I| ≤ m+1 (Carathéodory-type); if |I| = 1 then r\_1² ≤ m+1 (BP-type). The constraint Σ r\_i² ≤ m+1 interpolates between these extremes and is genuinely new for |I| > 1 with some r\_i > 1.

  Type: claim, Inference: assumption, Status: validated, Taint: unresolved

\item \textbf{Node 1.7.4:} (Novelty of the mixed bound). By Node 1.7.3, for an extreme point of the state space of A = direct\_sum M\_\{d\_i\}(C) under m affine constraints, the bound Sigma\_\{i in I\} r\_i\textasciicircum{}2 <= m+1 holds. We instantiate with A = M\_2(C) + M\_2(C) and m = 4: the bound gives Sigma r\_i\textasciicircum{}2 <= 5. The complete list of feasible rank pairs (r\_1, r\_2) with r\_i in \{0,1,2\} (where r\_i = 0 means block i is inactive) satisfying Sigma r\_i\textasciicircum{}2 <= 5 is: (r\_1, r\_2) in \{(0,1), (0,2), (1,0), (1,1), (1,2), (2,0), (2,1)\}. In particular, (r\_1, r\_2) = (2,2) is ruled out since 2\textasciicircum{}2 + 2\textasciicircum{}2 = 8 > 5. Thus with 4 constraints, an extreme point cannot have full rank in both blocks simultaneously. Novelty argument: Applying classical Barvinok-Pataki to each block independently means projecting the m = 4 constraints onto each M\_2(C) summand, giving at most 4 constraints per block, so BP per block yields r\_i\textasciicircum{}2 <= 5 for each i separately. This allows (2,2) since 4 <= 5 for each block. Hence the cross-block bound Sigma r\_i\textasciicircum{}2 <= 5 is strictly stronger than per-block BP. Remark on the M\_4(C) embedding: one could embed A = M\_2(C) + M\_2(C) into M\_4(C) via block-diagonal inclusion and apply classical BP to M\_4(C). This gives (r\_1+r\_2)\textasciicircum{}2 <= 5, i.e. r\_1+r\_2 <= 2, which is sometimes stronger (e.g. rules out (1,2)) but is not comparable in general to Sigma r\_i\textasciicircum{}2 <= 5 (e.g. (2,1) satisfies r\_1+r\_2 = 3 > 2 but Sigma r\_i\textasciicircum{}2 = 5 <= 5). The two bounds capture complementary structural information; the mixed bound Sigma r\_i\textasciicircum{}2 <= m+1 is thus genuinely new. [Inference: universal\_instantiation. Dependencies: 1.7.3.]

  Type: claim, Inference: assumption, Status: pending, Taint: unresolved

\end{enumerate}
\item \textbf{Node 1.8:} (Recovery: Infinite-Dimensional Case → Weis-Shirokov). When A = B(H) for separable H, the main theorem recovers Weis-Shirokov Theorem 3: extreme points of energy-constrained density operator sets \{ρ : Tr(H\_j ρ) ≤ E\_j, j=1,...,m\} have rank r satisfying r² ≤ m+1. For m ≤ 2, this forces r = 1 (pure states). Our framework extends this to states on non-Type-I algebras, where the face structure is richer.

  Type: claim, Inference: assumption, Status: pending, Taint: unresolved

\begin{enumerate}
\item \textbf{Node 1.8.1:} (Infinite-dimensional setup). Let H be a separable infinite-dimensional Hilbert space and A = B(H) the algebra of bounded operators. The predual of B(H) is the space T\_1(H) of trace-class operators (Takesaki, Theory of Operator Algebras I, §III.3), and the normal (= sigma-weakly continuous) states on B(H) are exactly the functionals phi\_rho(a) = Tr(rho a) where rho is a density operator (positive trace-class, Tr(rho) = 1). We denote the normal state space by S\_n(B(H)).

Topological distinctions: The full state space S(B(H)) subset B(H)* carries the weak-* topology on B(H)*. The normal state space S\_n(B(H)) carries the sigma-weak topology (equivalently, the topology induced by the trace-norm on T\_1(H)). Crucially, S\_n(B(H)) is NOT weak-* closed in B(H)* --- its weak-* closure includes singular states. However, S\_n(B(H)) is complete in the trace-norm topology and the set of density operators with Tr(rho) = 1 is trace-norm closed.

Normal states suffice for B(H): Since B(H) is a Type I factor, it has a unique irreducible representation (up to unitary equivalence), namely the identity representation on H. Consequently, every pure state (= extreme point of S(B(H))) is a vector state omega\_xi(a) = <xi, a xi>, and in particular is normal. More precisely, every state dominated by a normal state is normal. It follows that: (i) all faces of S(B(H)) generated by normal states consist entirely of normal states, and (ii) for any constraint set C = \{phi in S(B(H)) : phi(a\_j) = beta\_j\} with a\_j in B(H)\textasciicircum{}\{sa\}, every extreme point of C is a pure state of B(H) restricted to C, hence normal. Thus the downstream analysis can work exclusively with the normal state space S\_n(B(H)) and density operators.

For applications to bounded observables a\_1, ..., a\_m in B(H)\textasciicircum{}\{sa\}, the main theorem (node 1.3) applies directly. For unbounded Hamiltonians H\_j (self-adjoint operators affiliated with B(H)), the constraints rho -> Tr(H\_j rho) are generalized affine maps in the sense of Weis-Shirokov, and one must invoke node 1.2 in its generalized form; see node 1.8.3 for details.

  Type: claim, Inference: assumption, Status: pending, Taint: unresolved

\item \textbf{Node 1.8.2:} (GNS for density operators on B(H)). For a normal state φ\_ρ with ρ of finite rank r, the GNS representation π\_ρ is unitarily equivalent to id\textasciicircum{}\{⊕r\} (r copies of the identity representation of B(H) on H). The commutant is π\_ρ(B(H))' ≅ M\_r(C), so dim\_R(commutant\textasciicircum{}\{sa\}) = r². For rank-∞ density operators (e.g., thermal states), the commutant is infinite-dimensional.

  Type: claim, Inference: assumption, Status: validated, Taint: unresolved

\item \textbf{Node 1.8.3:} (Energy constraints and Weis-Shirokov recovery). Let H be a separable infinite-dimensional Hilbert space, A = B(H), and let S\_n(B(H)) denote the normal state space (identified with density operators via node 1.8.1). Consider m energy constraints defining C = \{rho in S\_n(B(H)) : Tr(H\_j rho) <= E\_j, j = 1,...,m\}.

CASE 1 (Bounded observables): Suppose H\_1, ..., H\_m in B(H)\textasciicircum{}\{sa\} are bounded self-adjoint operators. Then phi -> phi(H\_j) = Tr(H\_j rho) is a continuous affine functional on S\_n(B(H)) (in the trace-norm topology). The set C is a closed convex subset of the trace-norm compact set S\_n(B(H)), hence C is compact and convex, and by the Krein-Milman theorem C = conv(ext(C)) provided C is nonempty. For inequality constraints, node 1.2 (Weis-Shirokov, Corollary 5) applies directly to the compact convex set K = S\_n(B(H)) with m continuous affine functionals f\_j(rho) = Tr(H\_j rho) and sublevel constraints f\_j <= E\_j: if rho is an extreme point of C, then dim(F\_K(rho)) <= m. (Note: node 1.2 covers both equality and inequality constraints, and sublevel sets of continuous affine functionals.) Combining with the face-commutant correspondence (node 1.1) which gives dim(F\_K(rho)) = dim\_R(pi\_rho(B(H))'\textasciicircum{}\{sa\}) - 1, and the GNS computation (node 1.8.2) which gives dim\_R(pi\_rho(B(H))'\textasciicircum{}\{sa\}) = r\textasciicircum{}2 for rank-r density operators, we obtain r\textasciicircum{}2 - 1 <= m, i.e., r\textasciicircum{}2 <= m + 1. Dependencies: nodes 1.1, 1.2, 1.3, 1.8.1, 1.8.2.

CASE 2 (Unbounded Hamiltonians): Suppose H\_j are unbounded positive self-adjoint operators (affiliated with B(H) but NOT elements of B(H)). Then Tr(H\_j rho) is defined for all density operators rho via Tr(H\_j rho) = sum\_n <e\_n, H\_j\textasciicircum{}\{1/2\} rho H\_j\textasciicircum{}\{1/2\} e\_n> in [0, +infty], taking the value +infty when rho\textasciicircum{}\{1/2\} dom(H\_j\textasciicircum{}\{1/2\}) is not dense enough. The map f\_j : rho -> Tr(H\_j rho) is a generalized affine map in the sense of Weis-Shirokov (Definition 1 of [WS2021]) on the convex set K = S\_n(B(H)):

(a) Affinity on dom(f\_j): For rho\_1, rho\_2 with Tr(H\_j rho\_i) < infty, the identity Tr(H\_j (lambda rho\_1 + (1-lambda) rho\_2)) = lambda Tr(H\_j rho\_1) + (1-lambda) Tr(H\_j rho\_2) follows from linearity of the trace.

(b) Lower semicontinuity: Let \{P\_n\} be the spectral projections of H\_j onto [0,n]. Then P\_n H\_j P\_n in B(H)\textasciicircum{}\{sa\} is bounded, and rho -> Tr(P\_n H\_j P\_n rho) is a continuous affine functional on (S\_n(B(H)), trace-norm). By the spectral theorem, Tr(H\_j rho) = sup\_n Tr(P\_n H\_j P\_n rho), a supremum of continuous affine functions, hence lower semicontinuous in the trace-norm topology.

Since each f\_j is a lower semicontinuous generalized affine map, the sublevel set C = \{rho : f\_j(rho) <= E\_j, j=1,...,m\} falls within the scope of node 1.2 (Weis-Shirokov Theorem 2 / Corollary 5), which applies to lsc generalized affine maps. The face dimension bound dim(F\_K(rho)) <= m for extreme points rho of C follows. The face-commutant correspondence (node 1.1) and GNS computation (node 1.8.2) then yield r\textasciicircum{}2 <= m + 1 exactly as in Case 1. Note: node 1.3 as currently stated applies to bounded elements a\_j in A\textasciicircum{}\{sa\} with equality constraints; for the unbounded inequality case, we bypass node 1.3 and combine nodes 1.1, 1.2, and 1.8.2 directly.

RECOVERY OF WEIS-SHIROKOV: This recovers the content of Weis-Shirokov [J. Convex Anal. 28(3), 2021], specifically: Theorem 2 (face dimension bound for one generalized affine constraint), Corollary 5 (extension to m constraints), and Corollary 7 (the rank bound r <= sqrt(m+1) for extreme points of energy-constrained quantum states). Our derivation makes the algebraic mechanism transparent: the rank bound is the conjunction of the abstract face-dimension bound (node 1.2) with the B(H)-specific commutant computation (node 1.8.2). Dependencies: nodes 1.1, 1.2, 1.8.1, 1.8.2.

  Type: claim, Inference: assumption, Status: pending, Taint: unresolved

\item \textbf{Node 1.8.4:} (Special case m <= 2: automatic purity). Let H be a separable infinite-dimensional Hilbert space, A = B(H). Consider C = \{rho in S\_n(B(H)) : Tr(H\_j rho) <= E\_j, j = 1,...,m\} with m <= 2 constraints (bounded or unbounded H\_j). Let rho be an extreme point of C.

Argument via rank bound (from node 1.8.3): By node 1.8.3, r\textasciicircum{}2 <= m + 1 <= 3 where r = rank(rho). Since r is a positive integer (rho is a density operator, hence nonzero, so r >= 1), the inequality r\textasciicircum{}2 <= 3 forces r = 1. (Indeed, r = 2 would give r\textasciicircum{}2 = 4 > 3, a contradiction.) Therefore rho is a rank-1 projection, i.e., rho = |psi><psi| for some unit vector psi, and the corresponding state is pure.

Argument via face dimension gap (from node 1.9): The face dimension spectrum of S\_n(B(H)) is \{0\} union \{d\textasciicircum{}2 - 1 : d >= 2\} union \{+infty\} = \{0, 3, 8, 15, 24, ...\} union \{+infty\}. In particular, the face dimension gap (definition: face\_dimension\_gap) includes \{1, 2\}: there are no faces of S\_n(B(H)) of dimension 1 or 2. (This is because faces of S\_n(B(H)) generated by a rank-r state have dimension r\textasciicircum{}2 - 1, by the face-commutant correspondence (node 1.1) and the GNS computation (node 1.8.2), and r\textasciicircum{}2 - 1 in \{1,2\} has no positive integer solution.) By node 1.2 (Weis-Shirokov, Corollary 5), for an extreme point rho of C under m <= 2 constraints: dim(F\_\{S\_n(B(H))\}(rho)) <= m <= 2. Since dim(F\_\{S\_n(B(H))\}(rho)) cannot equal 1 or 2 (by the face dimension gap), and it is a nonneg integer, we must have dim(F\_\{S\_n(B(H))\}(rho)) = 0, meaning rho is an extreme point of S\_n(B(H)) itself, i.e., rho is a pure state.

This is the main quantum application of the Weis-Shirokov framework, corresponding to Corollary 7 of [Weis-Shirokov, J. Convex Anal. 28(3), 2021]. The face dimension gap argument shows that the purity conclusion is strictly stronger than what the rank bound alone gives: it provides a structural explanation (face gap) rather than a numerical coincidence. Dependencies: nodes 1.1, 1.2, 1.8.1, 1.8.2, 1.8.3, 1.9. Definitions used: face\_dimension\_gap, face\_commutant\_correspondence, Weis\_Shirokov\_theorem.

  Type: claim, Inference: assumption, Status: pending, Taint: unresolved

\item \textbf{Node 1.8.5:} (Extension beyond Type I). Our framework applies to states on arbitrary C*-algebras, not just B(H). For non-Type-I algebras (e.g., the hyperfinite II\_1 factor R, group C*-algebras of non-Type-I groups, free group factors L(F\_n)), the face structure of S(A) is much richer: there exist faces of all dimensions, the face dimension gap may differ, and the commutant need not be a direct sum of matrix algebras in general. However, the bound dim\_R(π\_φ(A)'\textasciicircum{}\{sa\}) ≤ m + 1 still forces the commutant to be finite-dimensional under finitely many constraints, and the Artin-Wedderburn decomposition applies to this finite-dimensional commutant regardless of the algebra type.

  Type: claim, Inference: assumption, Status: validated, Taint: unresolved

\end{enumerate}
\item \textbf{Node 1.9:} (Face Dimension Gap and Automatic Purity). Let A be a unital *-algebra. Define the face dimension gap of S(A) as the set G(A) = \{d ∈ N\_\{>0\} : no face of S(A) has dimension d\}. If G(A) ⊇ \{1, 2, ..., ℓ-1\} (i.e., the smallest positive face dimension is ≥ ℓ), then for any C = \{φ ∈ S(A) : φ(a\_j) = β\_j, j=1,...,m\} with m ≤ ℓ-1, every extreme point of C is automatically a pure state (extreme in S(A) itself). For B(H): G = \{1,2\}, so m ≤ 2 constraints suffice for automatic purity.

  Type: claim, Inference: assumption, Status: pending, Taint: unresolved

\begin{enumerate}
\item \textbf{Node 1.9.1:} (Definition of face dimension spectrum and gap). Let A be a unital C*-algebra (or more generally an Archimedean order unit space), so that the state space S(A) is weak-* compact. Define the face dimension spectrum Spec(A) = \{dim(F\_\{S(A)\}(phi)) : phi in S(A)\} subset of N\_0 union \{infty\}, where dim denotes the affine dimension of the face (equivalently, topological dimension; these agree for faces of compact convex sets, and both give infinity for infinite-dimensional faces). Note 0 in Spec(A) always: pure states exist by the Krein-Milman theorem (applicable since S(A) is compact convex). The face dimension gap is G(A) = \{d in N\_\{>0\} : d not in Spec(A)\}. Equivalently, by the face-commutant correspondence (node 1.1), Spec(A) = \{dim\_R(pi\_phi(A)'\textasciicircum{}\{sa\}) - 1 : phi in S(A)\} where pi\_phi is the GNS (cyclic) representation associated to phi. Thus d in Spec(A) iff there exists a GNS representation of A whose commutant has self-adjoint real dimension d + 1.

  Type: claim, Inference: assumption, Status: pending, Taint: unresolved

\item \textbf{Node 1.9.2:} (Computing the gap for key algebras — CORRECTED). (i) A = C(X) abelian, X infinite: Spec = \{0, 1, 2, ...\}, G = ∅ (measures supported on k points give face dim = k - 1). If |X| = n finite, Spec = \{0,...,n-1\}. (ii) A = M\_n(C): Spec = \{d²-1 : 1 ≤ d ≤ n\} = \{0, 3, 8, 15, ...\}, G ⊇ \{1, 2\}. (iii) A = B(H) inf-dim: Spec = \{0, 3, 8, 15, ...\} ∪ \{∞\}, G ⊇ \{1, 2\}. (iv) A = M\_2(C) ⊕ M\_2(C): CORRECTED — states can be supported on both blocks. Complete spectrum Spec = \{0, 1, 3, 4, 7\}, computed by enumerating all (|I|, r\_i) combinations: (|I|=1, r=1) dim 0; (|I|=1, r=2) dim 3; (|I|=2, r\_1=r\_2=1) dim 1; (|I|=2, one rank 1 one rank 2) dim 4; (|I|=2, r\_1=r\_2=2) dim 7. Correct gap G = \{2, 5, 6\}, NOT \{1, 2\}. The gap \{1,2\} only holds for SIMPLE algebras (N=1, d>=2). (v) Group C*-algebra C*(Gamma): Spec depends on dual; varies.

  Type: claim, Inference: assumption, Status: pending, Taint: unresolved

\item \textbf{Node 1.9.3:} (Automatic purity theorem). [Depends on: node 1.9.1 (definition of G(A) and Spec(A)), node 1.2 (Weis-Shirokov face-dimension bound), node 1.1 (face-commutant correspondence).] If the face dimension gap (Definition 1.9.1) satisfies G(A) contains \{1, 2, ..., l-1\} (i.e., the smallest positive face dimension in S(A) is >= l), then for any set C of states defined by m <= l-1 affine constraints, every extreme point of C is a pure state (extreme in S(A) itself). Proof: Let phi be extreme in C. By the Weis-Shirokov face-dimension bound (node 1.2), dim(F\_\{S(A)\}(phi)) <= m <= l-1. But dim(F\_\{S(A)\}(phi)) in N\_0 (face dimension is a non-negative integer: by the face-commutant correspondence (node 1.1), dim(F\_\{S(A)\}(phi)) = dim\_R(pi\_phi(A)'\textasciicircum{}\{sa\}) - 1, and the commutant's self-adjoint part is a finite-dimensional real vector space when the face is finite-dimensional). The gap hypothesis gives dim(F\_\{S(A)\}(phi)) not in \{1, ..., l-1\}. Combined with dim(F\_\{S(A)\}(phi)) <= l-1 and dim(F\_\{S(A)\}(phi)) in N\_0, the only possibility is dim(F\_\{S(A)\}(phi)) = 0, i.e., phi is pure.

  Type: claim, Inference: assumption, Status: pending, Taint: unresolved

\item \textbf{Node 1.9.4:} (Application of automatic purity to B(H)). [Depends on: node 1.9.1 (definition of G(A)), node 1.9.2 (spectrum computations), node 1.9.3 (automatic purity theorem).] For B(H) with dim(H) >= 2: the smallest positive face dimension is 3. For normal states (density operators), this arises from rank-2 states whose GNS commutant is M\_2(C), giving face dimension 2\textasciicircum{}2 - 1 = 3. For non-normal states (which exist when H is infinite-dimensional), the GNS representation factors through a quotient of B(H); such representations yield commutants that are either trivial (face dimension 0, i.e., pure) or infinite-dimensional (face dimension infinity), never dimension 1 or 2. Hence G(B(H)) contains \{1, 2\}, so l = 3. By the automatic purity theorem (node 1.9.3): under m <= 2 affine constraints, all extreme points of the constrained state set are pure. This recovers the specialization of Weis-Shirokov Corollary 7 to the quantum state space S(B(H)). Contrast with the abelian case: for C(X) with X infinite, Spec = \{0, 1, 2, ...\} (node 1.9.2(i)), G = empty, l = 1, so even 1 constraint can produce non-pure extreme points. For a concrete example, the set \{mu in P([0,1]) : integral of x dmu = 1/2\} has the extreme point (1/2)delta\_0 + (1/2)delta\_1, which is non-pure.

  Type: claim, Inference: assumption, Status: pending, Taint: unresolved

\item \textbf{Node 1.9.5:} (Quantitative gap for finite-dimensional C*-algebras). [Depends on: node 1.9.1 (definition of Spec(A) and G(A)), node 1.1 (face-commutant correspondence).] For A = direct\_sum\_\{i=1\}\textasciicircum{}N M\_\{d\_i\}(C) with all d\_i >= 2, the face dimension spectrum is Spec(A) = \{(sum\_\{i in I\} r\_i\textasciicircum{}2) - 1 : empty set \textbackslash\{\}!= I subseteq \{1,...,N\}, 1 <= r\_i <= d\_i for each i in I\} union \{0\}. This follows from the face-commutant correspondence (node 1.1): a state phi supported on blocks I with rank r\_i on the i-th block has commutant pi\_phi(A)' isomorphic to direct\_sum\_\{i in I\} M\_\{r\_i\}(C), with dim\_R(pi\_phi(A)'\textasciicircum{}\{sa\}) = sum\_\{i in I\} r\_i\textasciicircum{}2, giving face dimension (sum r\_i\textasciicircum{}2) - 1. The gap structure depends on N as follows. Case (a) N = 1 (simple algebra M\_d(C), d >= 2): Rank 1 gives face dim 0 (pure); rank 2 gives face dim 3; rank d gives face dim d\textasciicircum{}2 - 1. The smallest positive face dimension is 3, so G(A) contains \{1, 2\}. Case (b) N = 2 (two simple blocks, both d\_i >= 2): Face dimension 1 is achievable via |I| = 2 with r\_1 = r\_2 = 1 (sum r\_i\textasciicircum{}2 = 2, dim = 1). Face dimension 2 requires sum r\_i\textasciicircum{}2 = 3, which needs either |I| = 1 with r\textasciicircum{}2 = 3 (impossible, r must be integer) or |I| = 2 with r\_1\textasciicircum{}2 + r\_2\textasciicircum{}2 = 3, only solvable by \{r\_1, r\_2\} = \{1, sqrt(2)\} (invalid). So 2 is not in Spec(A), and G(A) contains \{2\} but not \{1\}. Case (c) N >= 3 (three or more simple blocks, all d\_i >= 2): Face dimension 1 is achievable (|I| = 2, r\_1 = r\_2 = 1). Face dimension 2 is also achievable (|I| = 3, r\_1 = r\_2 = r\_3 = 1, sum = 3, dim = 2). So G(A) contains neither 1 nor 2. Summary: (a) N = 1, d >= 2: G contains \{1, 2\}, automatic purity under m <= 2 constraints; (b) N = 2, all d\_i >= 2: G contains \{2\} but not \{1\}, automatic purity under m <= 1 constraint only; (c) N >= 3, all d\_i >= 2: G may not contain 1 or 2, no automatic purity guarantee from the gap alone. This is consistent with the corrected node 1.9.2(iv): Spec(M\_2(C) + M\_2(C)) = \{0, 1, 3, 4, 7\}, gap G = \{2, 5, 6\}.

  Type: claim, Inference: assumption, Status: pending, Taint: unresolved

\end{enumerate}
\item \textbf{Node 1.10:} (Application: Noncommutative Polynomial Optimization). At any finite level of the NPA hierarchy or state polynomial hierarchy, extremal solutions have GNS representations with commutant dimension Σ n\_j² ≤ m+1, where m is the number of constraints at that level. This implies the optimal state factors through a representation of bounded complexity.

  Type: claim, Inference: assumption, Status: pending, Taint: unresolved

\begin{enumerate}
\item \textbf{Node 1.10.1:} (NPA hierarchy as constrained state space — two regimes). We distinguish two regimes for applying the main theorem to the NPA hierarchy.

(A) FINITE LEVEL k (SDP spectrahedron). At level k, the NPA hierarchy defines a feasible set F\_k consisting of N\_k x N\_k positive semidefinite Hermitian matrices Gamma (moment matrices) satisfying m\_k independent affine scalar equations on the entries of Gamma. Here N\_k = |S\_k| is the number of words of length <= k in the operator variables, and the affine constraints encode: (a) commutation relations Gamma\_\{ij\} = Gamma\_\{kl\} for index pairs corresponding to words that agree modulo [A\_x\textasciicircum{}a, B\_y\textasciicircum{}b] = 0; (b) projector relations Gamma\_\{ij\} = 0 or Gamma\_\{ij\} = Gamma\_\{kl\} from E\_x\textasciicircum{}a E\_x\textasciicircum{}\{a'\} = delta\_\{aa'\} E\_x\textasciicircum{}a; (c) normalisation sum\_a E\_x\textasciicircum{}a = I. The number m\_k is the total count of independent such scalar affine equations on the upper-triangular (free) entries of Gamma, EXCLUDING the trace normalisation Tr(Gamma) = 1 (which is already part of the state space definition). Thus F\_k is a spectrahedron: a slice of the PSD cone H\_\{N\_k\}\textasciicircum{}+ by m\_k affine hyperplanes plus the trace-1 normalisation. This is precisely the setting of Node 1.6 with A = M\_\{N\_k\}(C) and m = m\_k affine constraints. The constraint Gamma >= 0 ensures the truncated linear functional phi is positive on sums of Hermitian squares of polynomials of degree <= k; this is STRICTLY WEAKER than positivity on the full *-algebra of the operators and is the reason the NPA hierarchy at finite k provides an outer RELAXATION of the quantum correlation set.

(B) CONVERGENCE (k -> infinity). In the limit k -> infinity, the hierarchy converges: the intersection of all F\_k equals the set of quantum correlations (Navascues-Pironio-Acin 2008). In this limit, feasible points correspond to states phi on the universal C*-algebra A of the Bell scenario (the C*-enveloping algebra of the free *-algebra on the operator variables modulo the measurement relations). This algebra A is infinite-dimensional and Archimedean (being a C*-algebra), so the main theorem (Node 1) applies directly to states on A subject to finitely many affine constraints. However, the number of affine constraints grows with k (m\_k -> infinity), so the bound Sigma n\_j\textasciicircum{}2 <= m\_k + 1 becomes vacuous in the limit unless one restricts to a finite number of constraints that are independent of k (e.g., fixing finitely many correlation entries).

This node is a corollary/application depending on Nodes 1 (main theorem), 1.3-1.4 (dimension bound and commutant decomposition), and 1.6 (matrix algebra recovery for SDP spectrahedra). It also depends on the definitions of NPA\_hierarchy and affinely\_constrained\_state\_space.

  Type: claim, Inference: assumption, Status: pending, Taint: unresolved

\item \textbf{Node 1.10.2:} (Structural consequence for extremal NPA solutions — corrected). By the main theorem (Node 1, applied via Node 1.10.1), any extreme point phi of the level-k feasible set F\_k (which is an SDP spectrahedron in M\_\{N\_k\}(C), cf. Node 1.10.1(A)) satisfies Sigma n\_j\textasciicircum{}2 <= m\_k + 1, where pi\_phi = direct\_sum sigma\_j\textasciicircum{}\{n\_j\} is the GNS decomposition of the state phi on M\_\{N\_k\}(C). This provides structural information about extreme points of the level-k SDP RELAXATION (not directly about quantum strategies, since at finite k the feasible set is an outer relaxation that includes pseudo-correlations not realizable by any quantum state or strategy):

(i) The commutant pi\_phi(M\_\{N\_k\}(C))' decomposes as direct\_sum M\_\{n\_j\}(C) with Sigma n\_j\textasciicircum{}2 <= m\_k + 1 (Artin-Wedderburn, Node 1.4). This bounds the COMMUTANT COMPLEXITY (number and multiplicities of isotypic components) but NOT the total Hilbert space dimension, since the irreducible dimensions dim(sigma\_j) are unconstrained by the bound. Concretely, the number of distinct irreducible summands k satisfies k <= m\_k + 1 (since n\_j >= 1 implies Sigma n\_j\textasciicircum{}2 >= k), and each multiplicity satisfies n\_j <= sqrt(m\_k + 1).

(ii) The commutant bound is the meaningful structural content: it constrains the multiplicity space dimensions, which determine how many independent ways the irreducible components can mix. For the NPA application at finite level k, since A = M\_\{N\_k\}(C) has a unique irreducible representation (the identity on C\textasciicircum{}\{N\_k\}), the decomposition simplifies to pi\_phi = id\textasciicircum{}\{r\} with r = rank(Gamma), and the bound becomes r\textasciicircum{}2 <= m\_k + 1, i.e., rank(Gamma) <= sqrt(m\_k + 1). This is the classical Barvinok-Pataki bound applied to the moment matrix.

(iii) Regarding SDP solvers: interior-point methods do NOT generically find extreme points. They converge to the ANALYTIC CENTER of the optimal face, which is typically NOT an extreme point (it is the point maximising a logarithmic barrier function over the optimal face). The analytic center has maximal rank among optimal solutions (Luo-Sturm-Zhang 2000). To obtain an extreme point (with the low-rank guarantee), one must apply a post-processing step: either (a) facial reduction to identify the minimal face, then a rank-minimisation heuristic, or (b) a randomised rounding procedure. Alternatively, first-order methods, augmented Lagrangian methods, or specialised low-rank solvers (Burer-Monteiro) may converge to low-rank solutions, but without guaranteed extremality.

This node is a corollary of the main theorem (Nodes 1.3-1.4) applied to the NPA setting (Node 1.10.1), not an assumption. It depends on Nodes 1, 1.3, 1.4, 1.6, and 1.10.1.

  Type: claim, Inference: assumption, Status: pending, Taint: unresolved

\item \textbf{Node 1.10.3:} (State polynomial optimisation — CORRECTED scope). The Klep-Magron-Volcic (KMV) state polynomial framework (Math. Program. 2023) generalises NPA by introducing 'state polynomials' — polynomials in noncommuting variables x AND a state symbol φ. The optimisation is over tuples (X, φ) where X are operators and φ a state; constraints involve state polynomials f(x, φ) that can be NONLINEAR in φ (e.g., φ(x²)φ(y) - φ(xy)² ≥ 0).

SCOPE RESTRICTION: Our main theorem (Nodes 1.3–1.4) applies to AFFINE constraints φ(a\_j) = β\_j or φ(a\_j) ≤ β\_j, where a\_j are fixed self-adjoint elements. This covers the SUBCASE of the KMV framework where all constraints are linear in φ (i.e., of the form φ(q\_j) ≤ β\_j for NC polynomials q\_j). This linear subcase is essentially the NPA framework (Node 1.10.1) and is already addressed there.

For the FULL KMV framework with nonlinear state polynomial constraints, our main theorem does not apply directly. However, two remarks connect our work to KMV:

(1) AT EACH FINITE LEVEL k, the KMV hierarchy produces an SDP relaxation (a spectrahedron defined by PSD and linear constraints on a moment matrix of size N\_k). The classical Barvinok-Pataki bound applies to this SDP, bounding the rank of the moment matrix: r² ≤ m\_k (complex case). Our theorem, applied to the truncated *-algebra A\_k = C⟨x\_1,...,x\_n⟩/J\_k (the quotient of the free *-algebra by the kernel at level k) with the m\_k linear constraints, gives the same rank bound with additional algebraic structure: the GNS decomposition Σ n\_j² ≤ m\_k + 1 reveals the multiplicity structure of extremal solutions, not just their rank.

(2) AT CONVERGENCE (k → ∞), if the hierarchy converges to the true optimum, the limiting feasible states live on a finitely-presented *-algebra A = C⟨x\_1,...,x\_n⟩/J with the original affine constraints. Our theorem applies to this limiting algebra when the constraints are affine in φ.

The genuine novelty beyond classical SDP rank bounds is structural: our theorem decomposes the commutant into Artin-Wedderburn blocks, revealing which irreducible representations appear and with what multiplicities, rather than merely bounding a single rank parameter. This is corollary/application depending on Nodes 1.3–1.4 and 1.10.1.

  Type: claim, Inference: assumption, Status: pending, Taint: unresolved

\item \textbf{Node 1.10.4:} (Quantum marginal problem application — CORRECTED). Let H = ⊗\_\{i ∈ I\} H\_i be a finite-dimensional multipartite Hilbert space with subsystems indexed by a finite set I, where dim(H\_i) = d\_i. Let S\_1, ..., S\_k ⊆ I be (not necessarily disjoint) subsets of subsystems, with corresponding marginal Hilbert spaces H\_\{S\_j\} = ⊗\_\{i ∈ S\_j\} H\_i of dimension D\_j = ∏\_\{i ∈ S\_j\} d\_i. Given target marginals ρ\_j ∈ S(B(H\_\{S\_j\})), the compatible state set is C = \{ρ ∈ S(B(H)) : Tr\_\{I∖S\_j\}(ρ) = ρ\_j, j = 1,...,k\}, where Tr\_\{I∖S\_j\} denotes the partial trace over all tensor factors not in S\_j.

To apply the main theorem (Nodes 1.3–1.4), we decompose each matrix marginal constraint into scalar affine constraints. Choose a Hermitian basis \{E\textasciicircum{}\{(j)\}\_α\}\_\{α=1\}\textasciicircum{}\{D\_j²\} for B(H\_\{S\_j\})\textasciicircum{}\{sa\}, e.g., the generalised Gell-Mann matrices plus (1/D\_j)·I. The constraint Tr\_\{I∖S\_j\}(ρ) = ρ\_j is equivalent to: Tr(ρ · (E\textasciicircum{}\{(j)\}\_α ⊗ I\_\{I∖S\_j\})) = Tr(ρ\_j · E\textasciicircum{}\{(j)\}\_α) for each α. Here a\textasciicircum{}\{(j)\}\_α := E\textasciicircum{}\{(j)\}\_α ⊗ I\_\{I∖S\_j\} ∈ B(H)\textasciicircum{}\{sa\} are the self-adjoint constraint elements.

CONSTRAINT COUNT (corrected): Each marginal gives D\_j² scalar constraints, but the α = identity component gives Tr(ρ · I\_H) = Tr(ρ\_j) = 1, which is automatically satisfied by the normalisation φ(1) = 1 in S(B(H)). So each marginal contributes D\_j² - 1 independent constraints. For NON-OVERLAPPING subsystems, m = Σ\_\{j=1\}\textasciicircum{}k (D\_j² - 1). For OVERLAPPING subsystems, additional linear dependencies arise (shared sub-marginals are determined by multiple constraints), so the actual number of independent constraints m is at most Σ\_\{j=1\}\textasciicircum{}k (D\_j² - 1) — possibly strictly less. The bound Σ n\_i² ≤ m + 1 remains valid with m being the number of independent affine constraints (since adding redundant constraints does not change the constrained set C).

For qubit marginals with non-overlapping single-qubit subsystems (D\_j = 2): each marginal contributes 2² - 1 = 3 independent constraints, so m = 3k and the bound is Σ n\_i² ≤ 3k + 1.

This node depends on the main theorem (Nodes 1.3–1.4) and is a corollary/application, not an assumption.

  Type: claim, Inference: assumption, Status: pending, Taint: unresolved

\end{enumerate}
\item \textbf{Node 1.11:} (Consistency Verification: Counterexample Check). The bound is consistent with the Weis-Shirokov counterexample (Bloch ball with 3 Pauli constraints): for A = M\_2(C) with m = 3, the bound gives r² ≤ 4, so r ≤ 2, and indeed the maximally mixed state (r = 2) is an extreme point of the constrained set, saturating the bound.

  Type: claim, Inference: assumption, Status: pending, Taint: unresolved

\begin{enumerate}
\item \textbf{Node 1.11.1:} (Counterexample setup). Consider A = M\_2(C) (2×2 matrices) with the three Pauli constraint observables a\_1 = σ\_x, a\_2 = σ\_y, a\_3 = σ\_z (the Pauli matrices, which form a basis for the traceless Hermitian 2×2 matrices). With m = 3 constraints, the constrained state space C = \{ρ ∈ S(M\_2(C)) : Tr(σ\_j ρ) = β\_j, j = 1,2,3\} consists of density matrices with prescribed Bloch vector (β\_1, β\_2, β\_3). For |β| < 1 (where β = (β\_1, β\_2, β\_3)), C = \{ρ\_β\} is a single point (the unique density matrix with that Bloch vector). For β = 0, this is the maximally mixed state ρ = I/2 of rank 2.

  Type: claim, Inference: assumption, Status: validated, Taint: unresolved

\item \textbf{Node 1.11.2:} (Checking the bound). The maximally mixed state ρ = I/2 has rank r = 2. Our bound gives r² = 4 ≤ m + 1 = 3 + 1 = 4. The bound is saturated (equality holds). This confirms the bound is tight in this case: with 3 constraints on M\_2(C), a rank-2 (non-pure) state can be extremal.

  Type: claim, Inference: assumption, Status: validated, Taint: unresolved

\item \textbf{Node 1.11.3:} (Why 3 constraints produce a non-pure extreme point). The state space S(M\_2(C)) is the Bloch ball (3-dimensional ball). The Bloch ball has dimension 3 = 2² - 1 = dim\_R(M\_2(C)\textasciicircum{}\{sa\}) - 1. Three affine constraints (fixing all 3 Bloch coordinates) determine a single point, which is trivially extreme. When this point is in the interior of the Bloch ball (i.e., β ≠ 0 but |β| < 1), it corresponds to a mixed state (rank 2). It is extreme in C (being the only point) but not extreme in S(M\_2(C)) (not on the boundary of the Bloch ball). This is consistent with the main theorem: m = 3 ≥ 3 = dim(face) = r² - 1, so non-pure extremal points are permitted.

  Type: claim, Inference: assumption, Status: pending, Taint: unresolved

\item \textbf{Node 1.11.4:} (Consistency with Im 2025). Im's Theorem 4.4 generalises Pataki's formula to general conic programs. For K = H\_2\textasciicircum{}+ (2×2 PSD Hermitian cone) and 3 affine constraints: dim(face(ρ, K)) - m + dim(correction) = 0 for extreme points. With face dimension r² = 4 and m = 3: the correction term dim(face(ρ, K)\textasciicircum{}⊥ ∩ R(A*)) accounts for the implicit redundancy. In this case, one of the 3 constraints is 'redundant' in the sense of being automatically satisfied at the face (the normalisation direction), yielding an effective m\_eff = 3 and face dim = 4 = 3 + 1, consistent with our bound. A full reconciliation with Im 2025 requires checking that the abstract version (via Weis-Shirokov) and the conic version (via Pataki 3.3.1) agree when applied to the same setting.

  Type: claim, Inference: assumption, Status: pending, Taint: unresolved

\end{enumerate}
\end{enumerate}

\end{document}
